\section{单位、分式域与环的特征}


\textbf{本节用到的记号：} $\mathbb Z/n\mathbb Z$ 是模 $n$ 的剩余类环；$R[x]$、$F[x]$ 是多项式环（第 1 节），$F$ 表示域；$\deg f$ 是多项式次数；$\binom pk$ 是二项式系数；$\hookrightarrow$ 表示单同态（嵌入），$\twoheadrightarrow$ 表示满同态；$\exists$ 是“存在”，$\forall$ 是“对所有”。

\subsection{单位与整环}


含幺环 $R$ 的单位群为
\[R^\times=\{u\in R:\exists v\in R,\ uv=vu=1\}.\]
它在乘法下构成群（单位就是“有乘法逆元的元素”）。若存在单位 $u\in R^\times$ 使 $r=us$，就说 $r$ 与 $s$ \textbf{相伴}，记作 $r\approx s$。域满足 $R^\times=R\setminus\{0\}$（每个非零元都可逆）；一般环中还可能有非零但不可逆的元素。

注意本节有两个相似的关系，符号特意分开：\textbf{相伴} $r\approx s$（用一个单位相乘联系起来）与后面分式域里\textbf{分数相等} $(a,b)\sim(c,d)$（交叉相乘相等）。相伴还有另一种写法：在整环中

\[
r\approx s\iff r\mid s\ \text{且}\ s\mid r\iff (r)=(s),
\]

即“互相整除”与“生成同一个主理想”是同一件事。这几种说法在整环中等价，在一般交换环中要小心，本讲只用整环情形。

还有一条常用常识：\textbf{零因子一定不是单位}。因为若 $ab=0$ 且 $a$ 可逆，则 $b=a^{-1}(ab)=0$。所以在 $\mathbb Z/n\mathbb Z$ 中每个元素要么是单位、要么是零因子，这就是 $\mathbb Z/n\mathbb Z$ 是域当且仅当 $n$ 为素数的另一个解释。

整环是交换含幺环、$1\ne0$，且没有非零零因子（与第 1 节的定义一致）。若 $ab=0$ 且 $a\ne0$，则 $b=0$；因此整环中约分成立。$\mathbb Z$ 与域上的多项式环 $F[x]$ 都是整环。

\subsubsection{$F[x]$ 的单位：一个必须写清前提的结论}


课堂上说“$F[x]$ 中的一个多项式可逆当且仅当它的常数项非零”。这句话在 \textbf{$F$ 是域} 时正确，但把它原样推广到任意系数环会出错。正确的分层是：

\begin{itemize}
\item \textbf{$F$ 是域}：$F[x]^\times=F\setminus\{0\}$，即单位恰为非零常数。所以 $1+x$ 不是单位——它的常数项虽然非零，但非常数。
\item \textbf{$R$ 是整环}：$R[x]^\times=R^\times$（[DF] §7.2 Proposition 4）。例如 $\mathbb Z[x]^\times=\{\pm1\}$，而 $2$ 在 $\mathbb Z[x]$ 中不是单位。
\item \textbf{$R$ 是一般交换含幺环}：$f=\sum_{i=0}^n a_ix^i$ 是 $R[x]$ 的单位，当且仅当 \textbf{$a_0$ 是 $R$ 的单位，且 $a_1,\dots,a_n$ 全是幂零元}。这里 $a$ 是\textbf{幂零元}指存在正整数 $m$ 使 $a^m=0$。这正是“常数项要可逆、其余部分要被幂零性吸收”的意思。
\end{itemize}


整环情形可以直接从次数证明：设 $f,g\ne0$ 且 $fg=1$。整环上首项系数之积非零，故

\[
\deg(fg)=\deg f+\deg g.
\]

左边次数为 $0$，所以 $\deg f=\deg g=0$：$f,g$ 都是 $R$ 中的元素，且互为逆元，于是 $f\in R^\times$。

一般环里的幂零条件在 $\mathbb Z/4\mathbb Z$ 上就能看到：$(1+2x)^2=1+4x+4x^2=1$，所以 $1+2x$ 是可逆的，尽管它不是常数。这里 $2$ 是 $\mathbb Z/4\mathbb Z$ 中的幂零元。做多项式环的题时，先确认系数环是哪一类，再决定用哪条结论。

\subsection{分式域}


设 $R$ 是整环。考虑 $R\times(R\setminus\{0\})$ 中的有序对 $(a,b)$，将
\[ (a,b)\sim(c,d)\iff ad=bc \]
视为等价关系，记等价类为 $a/b$。定义
\[\frac ab+\frac cd=\frac{ad+bc}{bd},\qquad \frac ab\frac cd=\frac{ac}{bd}.\]
得到一个域 $\operatorname{Frac}(R)$，其中 $a/b=0$ 当且仅当 $a=0$，且非零元 $a/b$ 的逆为 $b/a$。自然嵌入 $R\hookrightarrow\operatorname{Frac}(R)$ 为 $a\mapsto a/1$。

典型例子是 $\operatorname{Frac}(\mathbb Z)=\mathbb Q$，以及 $\operatorname{Frac}(F[x])=F(x)$，即有理函数域。

\textbf{泛性质（最小性）。} 上面的构造不只是“造出一个域”，它造出的是\textbf{包含 $R$ 的最小域}。精确地说：若 $K$ 是域、$\varphi:R\to K$ 是单同态（即把 $R$ 嵌入 $K$ 的一个同构副本），则存在唯一的域同态 $\widetilde{\varphi}:\operatorname{Frac}(R)\to K$ 使它在 $R$ 上与 $\varphi$ 一致：

\[
\widetilde{\varphi}(a/1)=\varphi(a)\quad(\forall a\in R).
\]

为什么唯一？因为 $\widetilde{\varphi}$ 是域同态，必须满足

\[
\widetilde{\varphi}(a/b)=\widetilde{\varphi}(a/1)\,\widetilde{\varphi}(1/b)=\varphi(a)\,\varphi(b)^{-1},
\]

右边完全由 $\varphi$ 决定，所以至多一种取法。存在性要把这个式子当作定义，再验证它与分数代表元的选取无关——这一步用的正是等价关系 $ad=bc$：若 $a/b=c/d$，即 $ad=bc$，两边施加 $\widetilde{\varphi}$ 即得 $\widetilde{\varphi}(a/b)=\widetilde{\varphi}(c/d)$。

\textbf{这是“最小域”的确切含义：} 任何包含 $R$ 的域 $K$ 都必须含有所有 $\varphi(a)\varphi(b)^{-1}$，也就是必须含有 $\operatorname{Frac}(R)$ 的一个同构副本；所以 $\operatorname{Frac}(R)$ 是包含 $R$ 的域中最小的一个。这也解释了 $\operatorname{Frac}(\mathbb Z)=\mathbb Q$ 的含义：$\mathbb Q$ 是包含 $\mathbb Z$ 的最小域，任何包含 $\mathbb Z$ 的域都含有 $\mathbb Q$。

\subsection{特征}


含幺环的特征定义为
\[\operatorname{char}(R)=\min\{n\ge1:n\cdot1_R=0\},\]
若不存在这样的正整数，则特征为 $0$。若 $R$ 是整环且特征非零，它必为素数：若 $n=ab$ 为合数，则 $(a1_R)(b1_R)=n1_R=0$；无零因子迫使 $a1_R=0$ 或 $b1_R=0$，与最小性矛盾。对一般含幺环，特征可以是合数，例如 $\mathbb Z/4\mathbb Z$ 的特征为 $4$。

特征为素数 $p$ 时有一个非常好用的推论：\textbf{若 $R$ 是特征 $p$ 的交换环，则对任意 $a,b\in R$}

\[
(a+b)^p=a^p+b^p .
\]

这是因为二项式展开 $(a+b)^p=\sum_{k=0}^p\binom pk a^kb^{p-k}$ 中，中间系数 $\binom pk$（$0<k<p$）都被 $p$ 整除，而它们在 $R$ 中作用为乘以 $p\cdot1_R=0$，于是中间项全部消失。有限域理论里的 Frobenius 映射 $x\mapsto x^p$ 就是这个等式。

对每个素数 $p$，$\mathbb F_p=\mathbb Z/p\mathbb Z$ 是含 $p$ 个元素的域，特征为 $p$。若 $R$ 是整环，则存在唯一的单位保持单同态

\[
\mathbb Z/\operatorname{char}(R)\mathbb Z\;\longrightarrow\;R,
\]

它在 $\operatorname{char}(R)=p$ 时给出 $\mathbb F_p\hookrightarrow R$，在 $\operatorname{char}(R)=0$ 时给出 $\mathbb Z\hookrightarrow R$。所以每个整环都含有一个“最小”的子环，称为\textbf{素子环}：特征为 $p$ 时是 $\mathbb F_p$，特征为 $0$ 时是 $\mathbb Z$。注意这里强调整环：一般含幺环上这个映射未必单射（例如 $\mathbb Z/4\mathbb Z$ 上的 $\mathbb Z\to\mathbb Z/4\mathbb Z$ 核为 $4\mathbb Z$）。

因此可以在 $\mathbb F_p[x]$ 上取多项式环，得到特征为 $p$ 的无限整环；再由分式域构造得到特征为 $p$ 的无限域 $\mathbb F_p(x)$。

\subsection{本讲检查清单}


1. \textbf{先确认约定}：本讲的环默认含幺且交换性随语境说明；谈“理想”时默认双边，谈“同态”时默认保幺。
2. \textbf{想证明某物是理想}，要验证非空、对减法封闭、并对\textbf{整个环}的左乘右乘封闭，而不是只对自身元素封闭。
3. \textbf{看到 $R/I$ 先确认 $I$ 是理想}，否则乘法根本不良定义。
4. \textbf{要证 $R/M$ 是域}，用理想对应定理把“$M$ 极大”翻译成“$R/M$ 没有非平凡理想”。
5. \textbf{要证 $R/P$ 是整环}，把“$ab\in P\Rightarrow a\in P$ 或 $b\in P$”翻译成“$R/P$ 无零因子”。
6. \textbf{中国剩余定理}的关键是 $I+J=R$ 给出满射；核总是 $I\cap J$，不需要互素。
7. \textbf{多项式环的单位}先看系数环是哪一类：域 → 非零常数；整环 → $R^\times$；一般环 → 常数项可逆 + 其余系数幂零。
8. \textbf{特征是否必为素数}只对整环成立；一般含幺环可以是合数。

\subsection{参考对照}


\begin{itemize}
\item [DF] §7.1--7.2（环、整环、单位与 $R[x]^\times=R^\times$）、§7.5（环的特征与素子环）、§9.3（分式域的一般构造）。
\item [DS] §10.1（整环与有限整环是域）、§10.3（特征与有限域）。
\item 原始材料：本地归档（手写原稿与录音转写），不纳入课程仓库。
\end{itemize}
